\section{Provider-arkitekturen}
\label{sec:providers}

Kjernen i provider-arkitekturen er grensesnittet \texttt{IVideoProvider},
som definerer åtte asynkrone metoder for romoperasjoner og
strømstyring. Ingen metoder har default-implementasjon. En provider som
ikke støtter en operasjon kaster \texttt{NotSupportedException}, som
\texttt{GlobalExceptionHandler} videre oversetter til HTTP-status \texttt{501 Not Implemented}. Provider velges
ved oppslag i ASP.NET Cores keyed DI-container via
\texttt{VideoProviderFactory}, som normaliserer nøkkelen til lowercase og
kaster \texttt{ArgumentException} ved ukjent provider. Grensesnittet er vist i kodeliste~\ref{lst:ivideoprovider}.

For at \texttt{VideoProviderFactory} skal kunne løse riktig provider ved kjøretid, registreres begge providers som keyed services i DI-containeren med en typet \texttt{HttpClient} og en \texttt{DelegatingHandler} for autentiseringshåndtering, som vist i kodeliste~\ref{lst:providerregistrering}.

\begin{lstlisting}[language={[Sharp]C}, label=lst:providerregistrering,
  caption={Registrering av providers som keyed services.}]
// HDO-VideoAPI/Extensions/ProviderExtensions.cs

public static IServiceCollection AddVideoProviders( ... )
{
    // AntMedia: Konfigurer og registrer provider med autentiseringshandler
    services.AddOptions<AntMediaConfig>()
        .Bind(config.GetSection("AntMedia"))
        // ...
    services.AddTransient<AntMediaAuthHandler>();
    services.AddHttpClient<AntMediaProvider>(client => { ... })
        .AddHttpMessageHandler<AntMediaAuthHandler>();
    // Registreres under nøkkel "antmedia" - muliggjør provider-valg i factory
    services.AddKeyedTransient<IVideoProvider>("antmedia", (sp, _) => sp.GetRequiredService<AntMediaProvider>());

    // NHN: Konfigurer og registrer provider med token-cache og autentiseringshandler
    services.AddOptions<NHNConfig>()
        .Bind(config.GetSection("NHN"))
        // ...
    services.AddSingleton<NHNTokenCache>();
    services.AddTransient<NHNAuthHandler>();
    services.AddHttpClient<NHNVideoProvider>(client => { ... })
        .AddHttpMessageHandler<NHNAuthHandler>();
    // Registreres under nøkkel "nhn" - muliggjør provider-valg i factory
    services.AddKeyedTransient<IVideoProvider>("nhn", (sp, _) => sp.GetRequiredService<NHNVideoProvider>());

    // Registrer factory som løser provider basert på nøkkel
    services.AddSingleton<IVideoProviderFactory, VideoProviderFactory>();
    return services;
}
\end{lstlisting}

% \texttt{AntMediaProvider} er strømorientert. Kjerneabstraksjonen i AntMedia er en broadcast med tilstandene \texttt{created}, \texttt{broadcasting} og \texttt{finished}. Start og stopp utføres som eksplisitte API-kall, og tilgang skjer via kortlevde WebRTC-tokens med 90 sekunders levetid. \texttt{AntMediaAuthHandler} signerer JWT-tokens lokalt med HMAC-SHA256 og cacher dem i fem minutter, med trådsikker oppdatering ved bruk av \texttt{SemaphoreSlim}.

% \texttt{NHNVideoProvider} er møteorientert. Et NHN-møte er en booking
% med start- og sluttid, og tilgang skjer via en Pexip-basert
% konferanse-URL med innebygd PIN som NHN genererer ved opprettelse. Siden
% møtet er aktivt så lenge bookingen eksisterer, finnes det ingen
% start- eller stopp-operasjon, og \texttt{SetRoomTokenAsync} er heller
% ikke aktuell. Disse tre metodene kaster \texttt{NotSupportedException},
% som kodeliste~\ref{lst:nhnsupported} viser.

Siden AntMedia er strømorientert (jf. delkapittel~\ref{sec:antmedia}), bygger \texttt{AntMediaProvider} rundt broadcast-abstraksjonen med tilstandene \texttt{created}, \texttt{broadcasting} og \texttt{finished}. Start og stopp utføres som eksplisitte API-kall, og tilgang skjer via kortlevde WebRTC-tokens med 90 sekunders levetid. \texttt{AntMediaAuthHandler} signerer JWT-tokens lokalt med HMAC-SHA256 og cacher dem i fem minutter, med trådsikker oppdatering via \texttt{SemaphoreSlim}.

Siden NHN Video er møteorientert (jf. delkapittel~\ref{sec:nhn}), modelleres et rom som en booking med start- og sluttid der tilgang skjer via en Pexip-basert konferanse-URL med innebygd PIN. Siden møtet er aktivt så lenge bookingen eksisterer, finnes det ingen start- eller stopp-operasjon, og \texttt{SetRoomTokenAsync} er heller ikke aktuell. Disse tre metodene kaster \texttt{NotSupportedException}, som kodeliste~\ref{lst:nhnsupported} viser.

\begin{lstlisting}[language={[Sharp]C}, label=lst:nhnsupported,
  caption={Operasjoner uten semantisk ekvivalent i NHN Video.}]
// HDO-VideoAPI/Providers/NHN/NHNVideoProvider.cs

public Task StartStreamAsync( ... ) =>
        throw new NotSupportedException("Start stream is not supported for NHN-rom.");

public Task StopStreamAsync( ... ) =>
        throw new NotSupportedException("Stop stream is not supported for NHN-rom.");

public Task SetRoomTokenAsync( ... ) =>
        throw new NotSupportedException("Set token is not supported for NHN-rom.");
\end{lstlisting}

% Autentisering mot NHN skjer via OAuth~2.0 password grant, og tokenet
% caches i singletonen \texttt{NHNTokenCache} i inntil 13 dager.
% Kontroller-laget ser ingen av disse forskjellene. En ny videomotor kan
% legges til ved å implementere \texttt{IVideoProvider} og registrere den
% med en ny nøkkel i \texttt{ProviderExtensions}.

Autentisering mot NHN skjer via OAuth~2.0 password grant, og tokenet caches i singletonen \texttt{NHNTokenCache} i inntil 13 dager. Controller-laget ser ingen av disse forskjellene — begge providers eksponerer det samme \texttt{IVideoProvider}-grensesnittet, og backend-spesifikk autentisering og domenehåndtering er fullstendig innkapslet i de respektive provider-implementasjonene.